\subsection{Candidate principles and primary sources}\label{sec:principios}
An AI agent (see \hyperref[sec:ai]{the declaration on the use of generative AI}) checked the available sources word for word against local copies: Chalmers (1995)~\cite{Chalmers1995}, Bayne (2010)~\cite{Bayne2010}, Schechter (2018)~\cite{Schechter2018}, three technical papers on integrated information~\cite{Kitazono2018,Oizumi2016PNAS,TokerSommer2019} and Mendola (2025)~\cite{Mendola2025}. The page numbers were taken from the running headers of plain-text copies and have not been checked against the printed editions. IIT~3.0~\cite{Oizumi2014}, whose postulates three rows of Table~\ref{tab:B-1} use, was not among these copies: those rows describe it without quotation, and the description has not been checked word for word against the original. Our answer, on the basis of these sources only, to the decisive question of Section~\ref{sec:mejor} (whether any available material rules out, with an independent foundation, that two original realizations with the same $\Phi_Q$ have different $u_{\rm orig}$) is \textbf{no}.

Table~\ref{tab:B-1} lists the candidate principles. The hypotheses OI$_\gamma$ and UT are stated in Definitions~\ref{hyp:OI} and~\ref{hyp:UT}, and Rob, B$_\psi^{\rm pres}$, N and the absence anchor in Definitions~\ref{hyp:Rob}, \ref{hyp:pres}, \ref{hyp:N} and~\ref{hyp:anchor}; the degree $u$ and the coexperience structure $\mathfrak U$ are those of Definition~\ref{def:unity}, and $B_\psi$ is the premise $u=\psi\circ\mathfrak U$.

\begin{table}[htbp]\centering\footnotesize
\begin{tabular}{>{\raggedright\arraybackslash}p{0.19\textwidth}>{\raggedright\arraybackslash}p{0.15\textwidth}>{\raggedright\arraybackslash}p{0.57\textwidth}}
\toprule
Principle & Scope & Foundation and limit\\
\midrule
OI$_\gamma$~\cite{Chalmers1995} & classes only & ``any system that has the same functional organization at a fine enough grain will have qualitatively identical conscious experiences'' (Chalmers, §1). It is independent, but it equates only isomorphic realizations, so it does not reach the cross-class residual (Theorem~\ref{thm:coverage}), which under Rob, with $\Phi_Q$ continuous, is a continuum in each interior fibre (Theorem~\ref{thm:continuum}, Proposition~\ref{prop:extension}).\\
Extension by fading & fibres, but (a) & A fading-type argument in the style of~\cite{Chalmers1995} infers that experience stays the same along a path of realizations from the fact that the subject's judgments and reports stay the same along it. Inside an isomorphism class the reports stay the same because the organization does. To extend the argument across the non-isomorphic realizations of a fibre of $\Phi_Q$, the reports about unity must stay the same across the fibre, that is, equal $\Phi_Q$ must give equal reports, which is $S_u$ at the level of reports. The extension therefore presupposes the bridge at the level of reports.\\
UT~\cite{Bayne2010,BayneChalmers2003} & fibres, but (b) & ``the simultaneous conscious states that S has at t will be subsumed by a single conscious state'' (Bayne, p.~16). Its domain is human beings (p.~17). With B$_\psi^{\rm pres}$ and N it gives $S_u$ in degenerate form: UT makes the coexperience structure total on the present aspects, B$_\psi^{\rm pres}$ makes $u$ a function of that structure, and N gives it the maximal value, so $u$ is constant on conscious realizations and $S_u$ holds there for every descriptor. Without N it is silent across different sets of present aspects: $u$ is then a function of the set of present aspects, and UT does not constrain that function.\\
IIT integration and exclusion~\cite{Oizumi2014} & (b) or (c) & Applied per complex, integration and exclusion make the experience of a complex a single one, so all its aspects are coexperienced, as UT says of a subject; with B$_\psi^{\rm pres}$ and N, $u$ then collapses as under UT. Applied per bearer, which may host several complexes, they make $u$ depend on how many complexes the bearer hosts, a count rather than a degree of coexperience, which changes the meaning of $u$. None of the sources checked argues that the degree of coexperience within a complex varies with $\Phi$.\\
IIT identity~\cite{Oizumi2014} & classes only & An experience is identical with the conceptual (cause--effect) structure that its complex specifies. The thesis equates $u$ only between realizations whose structures are isomorphic, and it entails $S_u$ for every compatible degree if and only if $\Phi_Q$ determines that structure up to isomorphism (Theorem~\ref{thm:impl}). Used to derive $S_u$, it thus needs a further premise that, with the thesis, already implies $S_u$: that $\Phi_Q$ determines the structure, or that $u$ depends on the structure only through $\Phi_Q$, which without a defense of its own is $S_u$ under another name.\\
IIT quantity; GNWT access & (c) & In IIT~3.0~\cite{Oizumi2014}, $\Phi^{\max}$ measures the quantity, or level, of consciousness of a complex; GNWT~\cite{DehaeneChangeux2011} concerns the global availability of information across brain systems. Coexperience is a relation among aspects within one experience (Definition~\ref{def:unity}). A level and an access relation are different targets, and taking either as $u$ changes the meaning of $u$.\\
Two unities (split brain)~\cite{Schechter2018}; Mendola~\cite{Mendola2025} & countermodels & ``not disunity or even simply partial unity but two unities'' (Schechter, p.~20). ``different sorts of phenomenal consciousness that cannot be organized along a single scale from least to greatest'' (Mendola, p.~82), said of forms and degrees of phenomenal consciousness, which Mendola carries over to unity in split-brain cases: ``a lesser degree or attenuated form of phenomenal consciousness and so unity that bridges the whole'' (p.~113). They count against scalar sufficiency, not for it.\\
\bottomrule
\end{tabular}
\caption{Candidate principles, with their scope, foundation and limit.}\label{tab:B-1}
\end{table}

In Table~\ref{tab:B-1}, (a) presupposes the bridge; (b) collapses $u$ into a constant or into a presence indicator; (c) changes the meaning of $u$. Together with the principles that reach only isomorphism classes, these are the ways out that Theorems~\ref{thm:impl}--\ref{thm:robusta} leave for the principles examined; changing the domain, as selection in the fibres does, is a further one (Section~\ref{sec:l3}).

\subsection{Best result and the premises that remain free}\label{sec:mejor}
\paragraph{Answer to the decisive question.} No available material offers an independent foundation for ruling out that two original realizations with the same $\Phi_Q$ have different $u_{\rm orig}$ when $u_{\rm orig}$ is not constant. By the theorems of Section~\ref{sec:intento}, the reason is structural:
\begin{enumerate}
\item The principles with a defense of their own that are aligned with $\Phi_Q$ (LSS$_\Delta$ and OI$_\gamma$, Definitions~\ref{hyp:LSS} and~\ref{hyp:OI}) equate only within isomorphism classes. The cross-class residual of each fibre remains intact (Theorems~\ref{thm:coverage} and~\ref{thm:residuo}), and under Rob, with $\Phi_Q$ continuous, it is a continuum (Theorem~\ref{thm:continuum}, Proposition~\ref{prop:extension}).
\item The transversal principles (coexperience structure, presence class, membership in a complex) imply $S_u$ for every degree compatible with them, that is, whatever function of the structure the degree is, only if $\Phi_Q$ determines that structure (Theorems~\ref{thm:impl} and~\ref{thm:robusta}). For one specified function $\psi$, $S_u$ requires only that structures that share a value of $\Phi_Q$ have equal values of $\psi$, which is condition~(ii) of SEG below with the level sets of $\psi$ as groups; it requires $\Phi_Q$ to separate two structures only when $\psi$ takes different values on them, and to determine the structures when $\psi$ takes different values on any two distinct structures, as an injective $\psi$ does. No $\psi$ has been specified (Definition~\ref{def:unity}), and for no nonconstant $\psi$ is that condition justified. If their co-occurrence with the levels of $\Phi_Q$ is connected, $S_u$ forces constancy (Theorem~\ref{thm:dicotomia}): the degree becomes the presence indicator or the normalization N, which are not adopted (Definitions~\ref{hyp:N} and~\ref{hyp:anchor}).
\end{enumerate}

\paragraph{SEG: non-degenerate $S_u$ in testable form.} Suppose that the degree is a function of the coexperience structure, $u=\psi\circ\mathfrak U$ ($B_\psi$, Definition~\ref{def:unity}). Then non-degenerate $S_u$, that is, $S_u$ with $u$ not constant on $D_0$, is equivalent to the following condition.

\begin{definition}[SEG]\label{def:SEG}
Assume $B_\psi$. The \emph{banded segregation} SEG holds on $D_0$ if the coexperience structures present in $D_0$, that is, the values $\mathfrak U(s)$ with $s\in D_0$, split into at least two groups such that:
\begin{enumerate}
\item[(i)] $\psi$ is constant on each group and takes different values on distinct groups;
\item[(ii)] no value of $\Phi_Q$ hosts structures from two distinct groups: if $s,t\in D_0$ and $\Phi_Q(s)=\Phi_Q(t)$, then $\mathfrak U(s)$ and $\mathfrak U(t)$ lie in the same group.
\end{enumerate}
\end{definition}

\begin{proposition}[\texttt{seg\_implies\_Su}, \texttt{segd\_iff}]\label{prop:seg}
Assume $B_\psi$. \textup{(a)} A split of the structures present in $D_0$ into at least two groups, with $\psi$ constant on each group and with condition~(ii) of Definition~\ref{def:SEG}, implies $S_u$. \textup{(b)} SEG holds if and only if $S_u$ holds and $u=\psi\circ\mathfrak U$ is not constant on $D_0$. \stL
\end{proposition}
\begin{proof}
(a) If $\Phi_Q(s)=\Phi_Q(t)$, condition~(ii) puts $\mathfrak U(s)$ and $\mathfrak U(t)$ in one group, on which $\psi$ is constant, so $u(s)=u(t)$. (b) If SEG holds, part~(a) gives $S_u$, and two realizations whose structures lie in distinct groups have different values of $u$, so $u$ is not constant. Conversely, if $S_u$ holds and $u$ is not constant on $D_0$, group the present structures by their value of $\psi$: condition~(i) holds by construction, there are at least two groups because $u$ is not constant, and condition~(ii) is $S_u$.
\end{proof}
In \texttt{FiltracionCociente.lean} (Appendix~\ref{sec:archivos}), the predicate \texttt{SEG} is the split of part~(a), without the requirement that $\psi$ differ between groups, and \texttt{SEGd} is Definition~\ref{def:SEG}.

SEG is not a new independent premise: it is non-degenerate $S_u$ written in the form in which it could be tested. Being equivalent is not being circular: a defense of SEG of its own, theoretical or empirical, would be a defense of non-degenerate $S_u$, and what is missing is such a defense, not another formulation. Its informative content is that the defended principles about coexperience ($B_\psi$, the presence class, membership in a complex) are transversal to $\Phi_Q$ (Section~\ref{sec:intento}), so $S_u$ requires $\Phi_Q$ to \emph{segregate} the coexperience structures into bands. The threshold law TL is the two-group case, presence versus absence. Under UT, B$_\psi^{\rm pres}$, N and the absence anchor, the degree takes one value on the conscious realizations and, by B$_\psi^{\rm pres}$ and the anchor, a smaller one on the non-conscious ones; SEG with these two groups says that no fibre of $\Phi_Q$ contains both a conscious and a non-conscious realization, and together with the order condition B (Definition~\ref{def:orderB}) it is TL (Definition~\ref{def:TL}, Theorem~\ref{thm:threshold}). Refuting SEG through a counterexample to $S_u$ would require three things:
\begin{itemize}
\item two original realizations with the same computed value of $\Phi_Q$;
\item coexperience structures determined by a calibrated observation channel (Q's observable-assignment rule A14--U is a stipulation, not a calibration);
\item a specified $\psi$ that distinguishes them.
\end{itemize}
None of the three is available: $\Phi_Q$ has not been selected (Definition~\ref{def:phiQ}), so none of its values has been computed; no calibrated channel is available; and no $\psi$ has been adopted (Definition~\ref{def:unity}). (By Proposition~\ref{prop:seg}(b), SEG also fails if $u=\psi\circ\mathfrak U$ is constant on $D_0$; that would not refute $S_u$.)

\paragraph{Free premises and decisions.} The following are open decisions in the specification of Q, not results to be proved: the selection of $\Phi_Q$, of the domain $D_0$ and of the grain $\gamma$ (Section~\ref{sec:enunciado}); and the comparison relation for the degree, a function $\psi$ of the coexperience structure or a substitute, together with the correspondence of aspects that it uses (Definition~\ref{def:unity}). The following remain unjustified for the original domain: LSS$_\Delta$, OI$_\gamma$, Rob, UT, $B_\psi$, B$_\psi^{\rm pres}$, N, the absence anchor, TL and the selection law SEL (Definition~\ref{def:SEL}). A9--U and A14--U are stipulated by Q, not derived (Definition~\ref{def:orderB}), and no result of this paper takes one of the decisions above or justifies one of these premises. The obligations of Q themselves remain open: $S_u$, or SEG in its testable form (Definition~\ref{def:SEG}), and the order condition B (Definition~\ref{def:orderB}). Table~\ref{tab:kinds} separates the kinds of item on which a derivation of the obligations of Q would rest.

\begin{table}[htbp]\centering\small
\begin{tabular}{>{\raggedright\arraybackslash}p{0.21\textwidth}>{\raggedright\arraybackslash}p{0.45\textwidth}>{\raggedright\arraybackslash}p{0.24\textwidth}}
\toprule
Kind & Items & Status in this paper\\
\midrule
Specification decisions (the author's) & $D_0$; the grain $\gamma$; $\Phi_Q$, with its boundary, units, state, dynamics, time scale, partitions, distance and selection conventions; the comparison relation of $u$ (a $\psi$, or a substitute) with its correspondences $f$; the reading of $\mathcal C$ & not taken; each result holds for every choice that satisfies its stated hypotheses\\
Stipulations of Q & A9--U ($S_u$ and B) and A14--U (Definition~\ref{def:orderB}) & stipulated, neither derived nor justified\\
Substantive premises & $B_\psi$, LSS$_\Delta$, OI$_\gamma$, Rob, UT, B$_\psi^{\rm pres}$, N, the absence anchor, TL and SEL on $D_0$; IPS$_K$ on the test bench of Section~\ref{sec:gauss} & examined; none adopted or justified\\
Results & theorems and propositions labelled \stL, \stM, \stA\ or \stS; numerical illustrations \stN & conditional on the hypotheses that each one states\\
\bottomrule
\end{tabular}
\caption{The kinds of item on which a derivation of the obligations of Q would rest; the reading of $\mathcal C$ enters only $S_q$, and IPS$_K$ only the test bench. No result of the last row takes a decision or justifies a premise of the rows above it.}\label{tab:kinds}
\end{table}

\paragraph{The other obligations of Q.} Q also carries sufficiency obligations for the components $c$, $q$, $\tau$ and $\sigma$, the order clause $\mathrm B$, the descent condition O, joint admissibility and conditions on processes (Definition~\ref{def:phiQ}). Section~\ref{sec:conditional} states which premises would discharge the five sufficiencies, $\mathrm B$ and O, and establishes none of them; apart from it, this paper treats only $S_u$. All these obligations remain open.

\subsection{What premises would discharge the obligations}\label{sec:conditional}
The following theorem applies the entailment criterion of Theorem~\ref{thm:impl} to each component, extends it to premises of order type, and adds the descent condition O. It changes neither the domain, nor the descriptor, nor the components, nor the obligations of Q, and it adopts none of the hypotheses of Section~\ref{sec:enunciado}.

\begin{theorem}[Conditional discharge of the sufficiency, order and descent obligations]\label{thm:conditional-discharge}
\stA\ Let $D_0$, $\Phi_Q$, the components $x\in\{c,q,u,\tau,\sigma\}$, their reduced inputs $d_x$ and the assignment $\mathbf e$ be as in Definition~\ref{def:phiQ}. For a relation $R$ on $D_0$, let $R^\ast$ be its reflexive and transitive closure and $\langle R\rangle$ the equivalence relation that it generates. A function $y$ on $D_0$ \emph{respects} $R$ if $s\,R\,t$ implies $y(s)=y(t)$, and a real-valued $y$ is \emph{nondecreasing along} $R$ if $s\,R\,t$ implies $y(s)\le y(t)$.
\begin{enumerate}
\item[(i)] \emph{Sufficiency, component by component.} Let $R_x$ be a relation on $D_0$. If $x$ respects $R_x$ and $\ker d_x\subseteq\langle R_x\rangle$, then $S_x$ holds. Conversely, let $X$ be a set with at least two elements; if every function $y:D_0\to X$ that respects $R_x$ satisfies $d_x(s)=d_x(t)\Rightarrow y(s)=y(t)$, then $\ker d_x\subseteq\langle R_x\rangle$. In particular, if $R_x$ is an equivalence relation contained in $\ker d_x$, every such $y$ satisfies that implication if and only if $R_x=\ker d_x$.
\item[(ii)] \emph{Unity through the coexperience structure.} Assume $B_\psi$ (Definition~\ref{def:unity}). Then $S_u$ is the inclusion $\ker\Phi_Q\subseteq\ker(\psi\circ\mathfrak U)$: realizations with equal $\Phi_Q$ may have different coexperience structures, provided that $\psi$ takes equal values on them. The relational sufficiency $\ker\Phi_Q\subseteq\ker\mathfrak U$ implies $S_u$ for every comparison $\psi$; it is not necessary in general for $S_u$ with one given $\psi$, even when $u$ is not constant; it is necessary for $S_u$ with one $\psi$ that takes different values on any two distinct structures occurring in one fibre of $\Phi_Q$, for instance an injective $\psi$ (for such a $\psi$, $S_u$ is equivalent to it), and for $S_u$ with every $\psi$ (Theorems~\ref{thm:impl} and~\ref{thm:robusta}). If $\mathfrak U$ respects a relation $R$ and $\ker\Phi_Q\subseteq\langle R\rangle$, then $\ker\Phi_Q\subseteq\ker\mathfrak U$.
\item[(iii)] \emph{Order.} Let $P$ be a relation on $D_0$, and put each real-valued function on $D_0$ that is nondecreasing along $P$ in place of $u$. (a) All these functions satisfy $S_u$ if and only if $\ker\Phi_Q\subseteq P^\ast\cap(P^\ast)^{-1}$. (b) All of them satisfy $\mathrm B$ if and only if $\Phi_Q(s)<\Phi_Q(t)$ implies $(s,t)\in P^\ast$. (c) All of them satisfy $S_u\wedge\mathrm B$ if and only if $\Phi_Q(s)\le\Phi_Q(t)$ implies $(s,t)\in P^\ast$. In particular, if $u$ is nondecreasing along $P$, each of these conditions implies the corresponding obligation. For $P=R\cup R^{-1}$ one has $P^\ast=\langle R\rangle$, a real-valued function is nondecreasing along $P$ exactly when it respects $R$, and (a) is Theorem~\ref{thm:impl}. If $S_u$ holds and $h_u$ is the function on $\Phi_Q(D_0)$ with $u=h_u\circ\Phi_Q$, then $\mathrm B$ holds if and only if $h_u$ is nondecreasing on $\Phi_Q(D_0)$.
\item[(iv)] \emph{Interventions.} Let $\mathcal T$ be a class of pairs of realizations of $D_0$ and $P_{\mathcal T}=\{(s,t)\in\mathcal T:\Phi_Q(s)<\Phi_Q(t)\}$, so that $\mathrm{IM}_{\mathcal T}$~\eqref{eq:im} says that $u$ is nondecreasing along $P_{\mathcal T}$. The real-valued functions on $D_0$ that are nondecreasing along $P_{\mathcal T}$ all satisfy $S_u$ if and only if $\Phi_Q$ is injective on $D_0$, and they all satisfy $\mathrm B$ if and only if any $s,t\in D_0$ with $\Phi_Q(s)<\Phi_Q(t)$ are joined by a chain of pairs of $P_{\mathcal T}$.
\item[(v)] \emph{Descent.} Let $\sim$ be a relation on $D_0$. If, for every $x$, $S_x$ holds and $s\sim t$ implies $d_x(s)=d_x(t)$, then $s\sim t$ implies $\mathbf e(s)=\mathbf e(t)$, which is condition O for $\sim$; when $\sim$ is an equivalence relation, this says that $\mathbf e$ factors through $D_0/{\sim}$. Conversely, if $s\sim t$ and $d_x(s)\ne d_x(t)$ for some $x$, then for every set $X$ with at least two elements some $y:D_0\to X$ satisfies $d_x(r)=d_x(r')\Rightarrow y(r)=y(r')$ for all $r,r'\in D_0$ and $y(s)\ne y(t)$.
\end{enumerate}
\end{theorem}
\begin{proof}
(i) If $d_x(s)=d_x(t)$ and $\ker d_x\subseteq\langle R_x\rangle$, there is a chain $s=r_0,\dots,r_n=t$ in which $r_{k-1}\,R_x\,r_k$ or $r_k\,R_x\,r_{k-1}$ for each $k$; $x$ takes one value along it, so $x(s)=x(t)$. Conversely, let $a\ne b$ in $X$ and $d_x(s)=d_x(t)$. The function $y$ with $y(r)=a$ if $r$ lies in the class of $t$ under $\langle R_x\rangle$ and $y(r)=b$ otherwise respects $R_x$, since $R_x$-related realizations lie in one class; hence $y(s)=y(t)=a$, and $s$ lies in the class of $t$. If $R_x$ is an equivalence relation contained in $\ker d_x$, then $\langle R_x\rangle=R_x\subseteq\ker d_x$, and the inclusion $\ker d_x\subseteq\langle R_x\rangle$ is the equality.

(ii) The first statement is $S_u$ for $u=\psi\circ\mathfrak U$, written out. If $\ker\Phi_Q\subseteq\ker\mathfrak U$, equal descriptors give equal structures, hence equal values of every $\psi$. For the failure of the converse in general, take a finite model, which is not a specification of Q: three realizations $s,t,r$ with descriptor values $0,0,1$ and three distinct structures on which $\psi=0,0,1$: for instance $E=\{a,b\}$ for the three episodes, $H_s=\{\{a\}\}$, $H_t=\{\{b\}\}$, $H_r=\{\{a\},\{b\},\{a,b\}\}$, and $\psi=1$ exactly when two aspects are coexperienced. In this model the sufficiency of the descriptor for $u=\psi\circ\mathfrak U$ holds, $u$ is not constant, and the descriptor does not determine $\mathfrak U$. The necessity for every $\psi$ is Theorem~\ref{thm:robusta}, or Theorem~\ref{thm:impl} with $R=\ker\mathfrak U$, whose indicators of single structures need no injective comparison; for one $\psi$ that separates the structures of a fibre, $\psi(\mathfrak U(s))=\psi(\mathfrak U(t))$ already gives $\mathfrak U(s)=\mathfrak U(t)$; and the last statement is the first part of~(i) with $\mathfrak U$ in place of $x$ and $\Phi_Q$ in place of $d_x$.

(iii) If $y$ is nondecreasing along $P$, induction on the length of a chain gives $y(s)\le y(t)$ whenever $(s,t)\in P^\ast$; this gives the ``if'' directions, in~(a) by antisymmetry. For the ``only if'' directions, fix $s$ and let $y_s$ be the indicator of $\{r:(s,r)\in P^\ast\}$. It is nondecreasing along $P$, since $(s,r)\in P^\ast$ and $r\,P\,r'$ give $(s,r')\in P^\ast$, and $y_s(s)=1$. In~(b), $\Phi_Q(s)<\Phi_Q(t)$ and $\mathrm B$ for $y_s$ give $1\le y_s(t)$, so $(s,t)\in P^\ast$; (c) is the same with $\le$ in place of $<$; in~(a), $S_u$ for $y_s$ gives $y_s(t)=1$, and $S_u$ for $y_t$ gives $(t,s)\in P^\ast$. If $P=R\cup R^{-1}$, then $P^\ast$ is symmetric, hence the smallest equivalence relation that contains $R$, and $y(s)\le y(t)\le y(s)$ on $R$-related pairs means $y(s)=y(t)$. Finally, if $S_u$ holds, $u=h_u\circ\Phi_Q$ with $h_u$ unique on $\Phi_Q(D_0)$ (Definition~\ref{def:phiQ}(e)); if $a<b$ lie in $\Phi_Q(D_0)$, with $a=\Phi_Q(s)$ and $b=\Phi_Q(t)$, $\mathrm B$ gives $h_u(a)=u(s)\le u(t)=h_u(b)$, and conversely a nondecreasing $h_u$ gives $\mathrm B$. This last equivalence is the reading of A9--U in Definition~\ref{def:orderB}(b): it restates $\mathrm B$ under $S_u$ and does not derive it.

(iv) Since $<$ is transitive, induction on chains shows that $(s,t)\in P_{\mathcal T}^\ast$ implies $s=t$ or $\Phi_Q(s)<\Phi_Q(t)$. Hence two distinct realizations with equal $\Phi_Q$ are never joined by a chain of pairs of $P_{\mathcal T}$, and the condition of~(iii)(a) holds exactly when every fibre of $\Phi_Q$ on $D_0$ has one element. The statement about $\mathrm B$ is~(iii)(b) with $P=P_{\mathcal T}$.

(v) If $s\sim t$, then $d_x(s)=d_x(t)$ for every $x$, so $S_x$ gives $x(s)=x(t)$ for every $x$, that is, $\mathbf e(s)=\mathbf e(t)$; a function that is constant on the classes of an equivalence relation factors through the quotient. Conversely, with $a\ne b$ in $X$, the function $y$ with $y(r)=a$ if $d_x(r)=d_x(s)$ and $y(r)=b$ otherwise is a function of $d_x$, and $y(s)=a\ne b=y(t)$.
\end{proof}
\emph{Formal check} \stL. The abstract form of Theorem~\ref{thm:conditional-discharge} is checked in Lean without Mathlib, in \path{CondicionalQ.lean} (namespace \texttt{CondicionalQ}); Table~\ref{tab:corr-1} lists the theorems that check each part. The realizations, the values of the components and of their reduced inputs are arbitrary types; the order of the descriptor values is a relation, irreflexive and transitive in~(iv), and in~(iii)(c) $\Phi_Q(s)\le\Phi_Q(t)$ is read as $\Phi_Q(s)<\Phi_Q(t)$ or $\Phi_Q(s)=\Phi_Q(t)$; and in~(iii) and~(iv) the degrees take values in a scale, a partial order with two values $\mathrm{lo}<\mathrm{hi}$ (\texttt{Scale}). The integers form a scale (\texttt{intScale}), and so do the real numbers, but that instance is not compiled, since Lean without Mathlib has no real numbers. In~(ii), \texttt{single\_psi\_example} takes the three structures as abstract codes, not as families of coexperienced sets, and the necessity for one $\psi$ that separates the structures of a fibre is not a separate compiled statement. Parts~(i), (iii)(b) and~(v) combine into one statement, \texttt{conditional\_discharge}; in a model on three realizations its bridge, coverage and order premises hold with $u$ not constant (\texttt{conditional\_discharge\_nonvacuous}).

\paragraph{What the theorem gives and what it does not.} Theorem~\ref{thm:conditional-discharge} is a conditional result. It separates premises of two kinds from its conclusions. The premises that $x$ respects $R_x$ and that $u$ is nondecreasing along $P$ are psychophysical hypotheses, and the construction of $\sim$ independently of Q is a physical task; the theorem justifies none of them. The inclusions of~(i), (iii) and~(v) are the second kind: when $R_x$, $P$ and $\sim$ are defined physically, they are conditions on $D_0$ and on the reduced inputs that become determinate propositions once $D_0$, $\gamma$ and the $d_x$ are selected, which Q has not done (Section~\ref{sec:laguna}, (L0)). Neither kind is an empirical validation, which would need a calibrated measurement channel for the components (Section~\ref{sec:mejor}). For $x=u$ and premises of the form $u(s)=u(t)$, part~(i) is Theorem~\ref{thm:impl} and adds nothing to Theorems~\ref{thm:coverage}--\ref{thm:robusta}: under Rob, with $\theta\mapsto\Phi_Q(s_\theta)$ continuous and $\Phi_Q(s_0)$ interior to its range, no relation contained in ${\cong_\gamma}$ covers $\ker\Phi_Q$ (Theorem~\ref{thm:coverage} and Proposition~\ref{prop:extension}), and with $R_u=\ker\Phi_Q$ the premise is $S_u$ itself. Part~(iv) makes the status of $\mathrm{IM}_{\mathcal T}$ (Section~\ref{sec:im}) exact: as a premise it entails $S_u$ for every compatible degree only when $\Phi_Q$ is injective on $D_0$, and it entails $\mathrm B$ only between realizations joined by interventions that raise $\Phi_Q$, so not between realizations of different bearers when $\mathcal T$ relates states of one bearer. The theorem does not treat joint admissibility or the process obligations, which Definition~\ref{def:phiQ} names without a formal statement; no conditional theorem about them can be stated before they are specified.

\paragraph{The quality component under one reading of $\mathcal C$ \stA\ \stL.} Q leaves open whether $\mathcal C$ in $d_q=(\Phi_Q,\mathcal C)$ is the unpointed wiring or the pointed isomorphism class at grain $\gamma$ (Section~\ref{sec:enunciado}). Under the second reading, Adm gives $\ker d_q={\cong_\gamma}$: $S_q$ then says exactly that $q$ is constant on the isomorphism classes at grain $\gamma$, the coverage condition of part~(i) holds with $R_q={\cong_\gamma}$, and $\Phi_Q$ plays no role in $S_q$ (\texttt{CondicionalQ.\allowbreak pointed\_reading\_kernel}, \texttt{CondicionalQ.\allowbreak Sq\_pointed}). OI$_\gamma$, which is not adopted here, yields that constancy on the classes contained in $D_{\rm cons}$ when $\gamma$ is fine enough in the sense of Definition~\ref{hyp:OI}, as stated after that definition for $u$. Under the first reading, $\ker d_q$ can contain pairs that are not isomorphic, and $S_q$ meets the coverage problem of $S_u$. Which reading Q intends is a decision of the author.
